% PROPOSED REVISED ABSTRACT
% File: ~/life/manuscript/sections/abstract.tex
%
% Goals of revision (per editorial-fit strategy):
%   1. Lead with clinical question, anchor on AIS epidemiology
%   2. Replace "energetic ground state" / "thermodynamic transition" /
%      "Bio-Gravitational" with biomechanical-control language that
%      a pediatric spine surgeon parses on first read.
%   3. Foreground the two strong correlations: 72% mechanosensor anisotropy
%      gap (p=0.011) and the derivative-gain proprioceptive-delay correlation
%      (r=0.983, p=2.7e-22).
%   4. Add the new molecular structural finding: mechanosensor proteins are
%      structurally rigid (low hinge density, p=0.003), supporting cooperative
%      mechanosensing.
%   5. End with falsifiable predictions, not metaphysical claims.

\begin{abstract}
\noindent\textbf{Background:} Adolescent Idiopathic Scoliosis (AIS) affects 2--4\% of adolescents, peaks at ages 11--15, and preferentially involves T8--T10 --- yet no unified mechanism explains the timing, anatomic localization, or 8:1 female predominance. We hypothesize that AIS is a predictable failure of an active, gravity-loaded postural-control system, not a stochastic genetic accident.

\noindent\textbf{Methods:} We model the developing spine as an information-elastic continuum: a Cosserat rod with stiffness modulated by HOX-patterned developmental information ($I(s)$) and stabilized by proprioceptive feedback against a gravitational load. Three independent in-silico strands: (i) cross-species allometric analysis across 12 vertebrate species (mouse to dolphin, $\mathcal{B}_g = EI/MgL^2$); (ii) AlphaFold-derived structural metrics across 61 candidate proteins (afcc pipeline); (iii) a delay-differential model of postural control through the adolescent growth spurt. \textbf{No human-subject Cobb measurements were used; all reported correlations are model-internal.}

\noindent\textbf{Results:} The S-curve emerges as the stable equilibrium configuration of the coupled information--mechanical system. Cross-species, only humans operate at $\mathcal{B}_g \approx 0.01$ where passive bending stiffness is insufficient and active maintenance is mandatory. AlphaFold analysis shows that a curated set of extended-filament/channel mechanosensors (Vimentin, Lamin~A/C, PIEZO2, SCRIB, PTK7 and related proteins) are 72\% more structurally elongated than metabolic regulators ($p=0.011$, Cohen's $d=1.19$, $n=23$); broader mechanotransduction-tagged sets dilute this signal (Results §3.4 sensitivity analysis). Mechanosensors are simultaneously more rigid in their interior ($0.93$ vs $3.12$ hinge regions per 100 residues, $p=0.003$, $n=61$) --- a signature of cooperative load-sensing in extended-filament/channel proteins. Within the model, simulated proprioceptive delay $\tau$ scales with body length while neural conduction velocity does not, creating a derivative-gain instability that closes the simulated postural control margin at the age window of AIS onset (in-silico correlation of $\Delta T$ with control-energy cost: $r=0.983$, $p=2.7\times 10^{-22}$). The simulated Cobb-angle response to varying tissue anisotropy shows strong protective scaling ($R^2=0.775$, $p<10^{-17}$, in-silico) with stability saturating at anisotropy $\geq 6$. We emphasize these are model-internal consistency checks; prospective external validation against longitudinal AIS cohort data is required.

\noindent\textbf{Conclusions:} AIS can be reframed from an isolated developmental defect to a predictable failure of an information-mechanical control system that maintains the S-curve against gravity. The framework yields three theoretical predictions, presented for hypothesis-generation only and not as clinical recommendations: (i) longitudinal proprioceptive latency should rise non-linearly through the growth spurt; (ii) under the model, mitochondrial cofactor availability (e.g., NAD$^+$) modulates predicted curve progression in simulation --- this generates a testable hypothesis but \textbf{does not endorse nutraceutical use in adolescents outside an IRB-approved trial}; (iii) ciliary length and cellular dithering should increase in scoliotic osteoblasts before macroscopic curve onset. Prospective testing in existing AIS cohorts under appropriate ethics oversight is the principal next step.
\end{abstract}
